
Deep learning models trained with empirical risk minimization can achieve high aggregate accuracy while failing on minority subgroups where spurious correlations do not hold. This work investigates whether per-sample loss trajectories during training encode information about spurious correlation vulnerability. Experiments on the Waterbirds benchmark show that trajectory features extracted from the first five epochs predict minority group membership with AUROC = 0.9452 $\pm$ 0.0072, exceeding the pre-specified threshold of 0.75. Initial loss ($L_{1})$ alone achieves AUROC = 0.9473, indicating that the discriminative signal is present from the first epoch. A controlled comparison of training regimes provides evidence for spurious-specificity: trajectory-based AUROC decreases by 0.29 under GroupDRO training but only by 0.01 under variance-matched random reweighting. A secondary hypothesis proposing delayed curvature stabilization in minority samples was not supported; curvature stabilized at epochs 2-3 for both groups. These results suggest that loss trajectory analysis may serve as a diagnostic tool for spurious correlation detection on this benchmark, though generalization to other datasets remains to be established.
